% Source: 21_Chapter_13_The_Sage_of_the_Ages_Augustine.txt
\bookchapter[The Sage of the Ages]{CHAPTER 13}{THE SAGE OF THE AGES}{Augustine}

\leadin{T}{HEY CALLED ROME THE} Eternal City. For 620 years, since the days of Hannibal, Rome had
seen no foreign invader outside its walls. Then suddenly in 410 Alaric, the Visigoth leader, with
his Arian hordes was besieging the city. The storm would soon break. Everyone knew that, but when?

The first peace party ventured beyond the wall to bargain with Alaric. They begged for mercy. They
asked for terms. All your gold, all your silver, came the reply. Above all, your German slaves! The
Romans continued to plead as the misery within the walls mounted.

Finally, the Visigoths charged through the gates and plundered the city, temple by temple, palace by
palace. Devastation and ruin were everywhere—except in the churches. When Alaric, proclaiming
himself a Christian, inspected the booty, he separated the church treasures from the rest and had
his soldiers carry the sacred vessels through the streets to the churches dedicated to Peter and to
Paul and leave them there.

A short time later the Visigoth and his troops withdrew from Rome, but the world would never be the
same. The glory of “the queen of cities” was gone. Eternal Rome was not eternal. “My voice sticks in
my throat,” said Jerome. “The City that took the whole world captive is itself taken.” He uttered
the sentiments of all, Christians and heathen.

Shocked Romans pointed to the spots where statues of the ancient gods had stood. They had made Rome
great. Perhaps they would have saved the city. Were they angry because the emperors had so recently
turned to the Christian God?

Refugees from Rome fled in all directions. In the little North African seaport of Hippo, a slim and
beardless figure, with a shaven head and sharp features, watched the arrival of the troubled
refugees. He heard their probing questions, their expressions of doubt. Someone should offer an
explanation for the ravaging of Rome. Then and there Aurelius Augustinus, the Bishop of Hippo,
decided to pursue the questions: Why had Rome fallen? Would the ruin of the Eternal City mean the
collapse of Christianity? Was the end of the world at hand?

Augustine’s answer to these questions not only provided light for the dark days just ahead, but a
philosophy for the foundations of Christendom. To this day Christians feel the impact of his mind
and soul. Roman Catholicism draws upon Augustine’s doctrine of the church and Protestants upon his
views of sin and grace.

\section{PREPARATION FOR GREATNESS}

The great North African was born on November 13, 354, in Tagaste, a little town in the hill country
of Numidia, a region we know as Algeria. His father, Patricius, was an easygoing heathen; his
mother, Monica, an eager Christian. In spite of their limited resources, the couple was determined
to give Augustine the best education available. Thus he went to school, first near home and then in
the North African capital, Carthage. In the capital he found sexual temptations irresistible. He
fell in love with a girl who gave him a son, Adeodatus. They lived together for thirteen years, but
Augustine always felt that sex was his defiling passion. It colored his conceptions of sin, and it
marked the depravity from which he later felt himself rescued by God’s grace.

His higher nature, however, constantly asserted itself. A treatise by Cicero, the Latin author, fell
into his hands in his nineteenth year and convinced him intellectually that he should make truth his
life’s search. The old temptations, however, still assailed him, and like Paul, he felt that two
warriors, a higher and a lower, were struggling in him for mastery. In his conflicts he turned to
the Bible, but it had no appeal to him. Its style seemed crude and barbarous to him.

For a time Augustine tried Manicheanism, a persecuted faith in the Roman Empire, but one peculiarly
appealing to a man of passion who felt two tendencies at war within him. Mani, its founder, had
taught in Persia and had met there a martyr’s death by crucifixion in 276 or 277. The fundamental
belief of the religion pictured the universe as the scene of an eternal conflict of two powers, the
one good and the other evil. Man, as we know him, is a mixed product: the spiritual part of his
nature consists of the good element, the physical of the evil. His task, therefore, is to free the
good in him from the evil, and this can be accomplished by prayer, but especially by abstinence from
all the enjoyments of evil: riches, lust, wine, meats, luxurious houses, and the like.

Manichaeism, like Gnosticism, taught that the true spiritual Jesus had no material body and did not
actually die. His purpose was to teach men the way out of the kingdom of darkness into the kingdom
of light. Like the Gnostics, the Manicheans held that much of the New Testament is true, but they
rejected everything in it that seemed to suggest Christ’s real sufferings, and they discarded the
Old Testament altogether.

Augustine remained an eager Manichean for nine years, from 374 to 383, before dissatisfaction with
its teachings arose in his mind. During these years he taught grammar in his native Tagaste, and
rhetoric in Carthage, and though inwardly doubting the truth of the Manichean philosophy, it was at
the suggestion of Manichean friends in 383 that he moved to Rome.

Shortly after his arrival in the capital he secured a professorship in the State University in Milan
(384) and moved to the northern city. His widowed mother and some of his African friends soon joined
him. He was now thirty years old, at the summit of a career, with dazzling prospects of success
before him. More than ever, however, he was deeply dissatisfied with his life. He callously
separated from his mistress, Adeodatus’s mother, to become engaged to a young woman of wealth and
position; but he could not master his passions. He found himself in “a whirl of vicious
love-making.” “There’s nothing so powerful,” he later observed, “in drawing the spirit of a man
downward as the caresses of a woman.” His inner conflicts became almost unbearable.

While living in Milan, however, Augustine came under the powerful preaching of Bishop Ambrose. He
went to church first to study Ambrose’s preaching style, but before long the message reached his
soul. In Ambrose he discovered that Christianity could be both eloquent and intelligent, and that
the troublesome stories in the Old Testament could be interpreted as allegories.

The final stimulus to Augustine’s conversion seems to have been the personal example of the monks.
When a friend told him the story of Anthony and the Egyptian hermits—how they withstood the
temptations of the world—Augustine felt a burning sense of shame. How could these unlearned men win
such spiritual victories when he, with all his education, knew only defeat? His sense of sin and his
powerlessness were profoundly stirred.

Matters came to a head as he walked through his garden in agony. He heard the singsong voice of a
child saying, “Take it and read it.” He picked up a New Testament. His eyes fell on the words
perfectly suited to his mood: “not in rioting and drunkenness, not in chambering and wantonness, not
in strife and envying. But put ye on the Lord Jesus Christ, and make not provision for the flesh, to
fulfil the lusts thereof” (Rom. 13:13–14, KJV). “Instantly,” said Augustine, “as I reached the end
of this sentence, it was as if the light of peace was poured into my heart, and all the shades of
doubt faded away.”

On the eve of the following Easter, 387, Augustine, with his son, Adeodatus, and his friend,
Alypius, was baptized by Ambrose in Milan. “The unrest of our past life,” he said, “receded from
us.” A few months later, accompanied by his mother, he set out for North Africa a different man. On
the journey, however, near Rome, his mother died. And during the autumn of 388, once again settled
in Tagaste, he lost his son, adding to the grief he already felt from the death of his mother.
Augustine was now as eager to leave the world as he had once been to plunge into it. It was not to
be. His gifts were too well-known and the need for leadership in the church was too great for him to
be left in peace.

Three years later at Hippo, by popular demand but against his will, Augustine was ordained a priest.
Soon, at the request of Bishop Valerius, he was chosen assistant bishop of the church, and a year
later, upon the death of Valerius, Augustine succeeded to the leadership of the church in Hippo. He
was forty-three years old and for the next thirty-three years, until his death in 430, he stood in
the center of the storms of his time.

\section{THE DONATIST CONTROVERSY}

North African Christianity was still torn by a passionate conflict between Catholics and a movement
called Donatism. The controversy was long-standing and deep-seated. A bishop of Hippo could scarcely
avoid speaking to the issue.

When Augustine stepped into church leadership, Donatism was almost a hundred years old. The movement
stood for a holy church, for church discipline, and for the unflinching resistance of unworthy
bishops. The Catholics, said the Donatists, had surrendered all of these by ordaining immoral
priests and bishops.

The Donatist name arose from Donatus, an early bishop of Carthage (313–355) who led the protest
against Catholic practices. Donatist charges centered on the fact that certain Catholic bishops had
handed over the Scriptures to be burned during the persecution under Diocletian. Such an act, the
Donatists insisted, was a serious sin of apostasy. Since the Catholic pastors were ordained by
bishops who had sinned so grievously, the Donatists believed they, rather than the Catholics,
constituted the true church of Christ. During Augustine’s time the Donatists were still widespread
in North Africa, and in some areas they constituted a majority.

Augustine rejected the Donatist’s view of a pure church. Until the day of judgment, he said, the
church must be a mixed multitude. Both good and bad people are in it. To support this idea he
appealed to Jesus’ parable of the wheat and tares (Matt. 13:24–30), overlooking the fact that Jesus
was not speaking of the church but of the whole world.

Augustine also set forth a different understanding of the sacraments. The Donatists argued that the
validity of the sacrament depends upon the moral standing of the minister. Augustine said no. The
sacrament does not belong to the minister but to Christ. The priest’s acts are really God’s because
he has placed the sacraments in the hands of the properly ordained minister. All that is required of
the priest is his awareness that he administers God’s grace for the whole church.

Such a view makes the priest the channel for grace to the members of the church. Thus Augustine
added his considerable influence to his priestly (sacerdotal) view of the church that reached such
unfortunate extremes in medieval Catholicism.

His defense of the Catholic Church in the Donatist controversy also led Augustine to support the use
of force in the suppression of the rivals. Initially he was strongly opposed to coercion. But step
by step he came to another view. In the face of Donatist resistance to the government’s mounting
pressure, he came to accept the use of force in a religious issue. What looks like harsh action, he
said, may bring the offender to recognize its justice. Had not the Lord himself in the parable said,
“Compel people to come in” (Luke 14:23)? Thus Augustine’s prestige was made available for those in
later ages who justified the ruthless acts of the Inquisition against Christian dissenters.

\section{OF SIN AND GRACE}

If the Donatist controversy called forth the churchman in Augustine, the Pelagian debate evoked his
emphases on salvation by grace.

Pelagius was a British monk who came to North Africa from Rome. A disciple who accompanied him,
Coelestius, had hopes of securing ordination as a priest in Carthage, but he found little support in
lands dominated by Augustine. As soon as Coelestius’s views appeared in Carthage, they were
repudiated.

As a result Coelestius headed for the East where Pelagius had gone. Churches there were more
receptive to Pelagius’s teaching. But Augustine launched a strenuous literary attack on Pelagianism.
By 419 the Pelagians were banished by the Emperor Honorius, and in 431 they were condemned by the
General Council of the Church meeting at Ephesus.

Why? What did Pelagius teach to arouse Augustine’s vigorous opposition? The monk denied that human
sin is inherited from Adam. Man, he said, is free to act righteously or sinfully. Moreover, death is
not a consequence of Adam’s disobedience. Adam, indeed, introduced sin into the world, but only by
his corrupting example. There is no direct connection between his sin and the moral condition of
mankind. Almost all the human race has sinned; but it is possible not to sin, and some people have
in fact lived without sin. God predestinates no one, except in the sense that he foresees who will
believe and who will reject his gracious influences. His forgiveness comes to all who exercise
“faith alone”; but, once forgiven, man has power of himself to live pleasing to God. Thus Pelagius
found no real need for the special enabling power of the Holy Spirit. His idea of the Christian life
was practically the Stoic conception of ascetic self-control.

All this was in sharp contrast to Augustine’s own experience. He sensed profoundly the depth of his
sin and hence the greatness of God’s salvation. He felt that nothing less than irresistible divine
power (grace) could have saved him from his sin and only constantly inflowing divine grace could
keep him in the Christian life. His Christian ideal was no Stoic self-control but love for
righteousness infused by the Spirit of God.

In Augustine’s view, Adam’s sin had enormous consequences. His power to do right was gone. In a
word, he died—spiritually, and soon physically. But he was not alone in his ruin. Augustine taught
that the whole human race was “in Adam” and shared his fall. Mankind became a “mass of corruption,”
incapable of any good (saving) act. Every individual, from earliest infancy to old age, deserves
nothing but damnation.

Since man of himself can do nothing good, all power to do good must be the free gift of God, that
is, grace. Out of the mass of the fallen race, God chooses some to receive this grace, which comes
to them from the work of Christ, and ordinarily through the church and especially through its
sacraments. All who receive baptism receive regenerating grace, which gives man back his freedom to
serve God, though that service is imperfect even at best, and is maintained only by the constant
gift of more grace.

As Williston Walker explains, “Those to whom God does not send his grace are lost. Nor can any man
be sure, even if he now enjoys God’s grace, that he will be saved. Only those to whom God gives the
added grace of perseverance, that is, who have divine aid to the end of life, will be redeemed.”
Man, therefore, has no power or worthiness of himself; his salvation is fully from God.

Augustine, like Paul, had been “apprehended” by the grace of God. He entered the Christian faith by
what seemed to him catastrophes. The great original catastrophe was his condition at birth, sinful
alienation from God. And the only freedom from that catastrophe was a new birth.

The Protestant Reformers would depend greatly upon Augustine’s view that mankind is incapacitated by
sin’s grip and the notion that only God’s grace could rehabilitate man. These teachings fit well
with Augustine’s teaching of predestination, which the Reformers expanded. They disregarded how
Augustine tied salvation and grace to the sacraments of the church.

Pelagius’s teaching was condemned by the great council of Ephesus in 431. He had failed to grasp
sin’s great distorting effect and how turning toward God without his grace was impossible. Yet over
the next century the church expressed independence from Augustine’s ideas of irresistible grace and
predestination. Some critics argued that Augustine had broken the longstanding practice of embracing
human freedom. Vincent of Lerins complained of this innovation when he wrote that a Christian should
believe what has been believed “always, everywhere, and by all.”

\section{THE CITY OF GOD ON THE RUINS OF ROME}

Augustine was fifty-six years old when he was told that Rome had been sacked. It must have been a
dramatic moment in his life. He welcomed the first refugees from Rome and set to work finding
quarters for them and encouraging his enlarged flock. In a sermon delivered at the time, he compared
the capture of Rome with the judgment of Sodom. There had been a great deal of destruction, but
cities, he said, consist of men, not walls. Unlike Sodom, Rome had been chastised but not destroyed.

Soon, he turned to the deeper questions of the relations between earthly cities like Rome, which
have their day, rising and falling like everything in time, and the Heavenly City or City of God,
which is everlasting. This question occupied him for sixteen years, almost to the end of his life,
and resulted in his great work, The City of God, which directly or indirectly influenced the thought
of Christians on what they owed to God and what to Caesar through the succeeding fifteen centuries.

From Adam to the end of time, Augustine wrote, humanity falls into two cities: the mass of the
godless, who live the life of earthly men; and the company of spiritual men, born of grace and
called to the City of God for all eternity. The Worldly City, said Augustine, is united by the
common love for temporal things. The City of God is bound together by the love of God.

What drove the Romans to their great achievements except the praise of men?

\begin{blockquote}
What else was there for them to love save glory? For, through glory, they desired to have a kind of
life after death on the lips of those who praised them~.~.~.~. The Heavenly City outshines Rome,
beyond comparison. There, instead of victory, is truth; instead of high rank, holiness; instead of
peace, felicity; instead of life, eternity~.~.~.~.
\end{blockquote}

And what about the church and state? Augustine considered the church the only human community that
worked for the building of the City of God. The state had its place in suppressing crime and
preserving peace, but since the state was based on the power of sin, it must submit to the laws of
the Christian church.

The grandeur of this spiritual vision made The City of God the most beloved of Augustine’s works
throughout the early Middle Ages. It gave a spiritual interpretation to the woes the world was
suffering. The present might be bad, but better things are to come. The golden age—the kingdom of
God—is in the future, not in the fading splendors of a worldly kingdom that could only crumble and
fall.

The older Augustine became, the more difficult his life became. In his seventy-sixth year the
barbarian Vandals had crossed the Straits of Gibraltar and were sweeping east toward Hippo. In his
closing days he had the penitential psalms copied on parchment and fixed to the wall of his room so
he could read them in bed. The dying man believed the end of the world to be at hand. On August 28,
430, Augustine found peace at last in Christ.

The Vandal siege of Hippo lasted fourteen months. When they finally breached the walls in August
431, they found most of the people dying or dead from hunger.

\readinghead{Suggestions for Further Reading}
\begin{readinglist}
\entry *Bonner, Gerald. St. Augustine of Hippo: Life and Controversies. London: SMC, 1963.
\entry Brown, Peter. Augustine of Hippo: A Biography. Berkeley: University of California Press,
  2000.
\entry Hansel, Robert R. The Life of Saint Augustine. New York: Franklin Watts, 1969.
\entry Knowles, Andrew and Pachomius Penkett. Augustine and His World. Downers Grove, IL:
  InterVarsity Press, 2004.
\entry Martin, Thomas F. Our Restless Heart: The Augustinian Tradition. Maryknoll, NY: Orbis, 2003.
\entry O’Meara, John J. The Young Augustine. Staten Island: Alba House, 1965.
\end{readinglist}
